\documentclass[11pt,a4paper,sans]{moderncv}        % possible options include font size ('10pt', '11pt' and '12pt'), paper size ('a4paper', 'letterpaper', 'a5paper', 'legalpaper', 'executivepaper' and 'landscape') and font family ('sans' and 'roman')

\usepackage[default,light,opentype]{sourcesanspro}
\usepackage{microtype}

% moderncv themes
\moderncvstyle{classic}                             % style options are 'casual' (default), 'classic', 'banking', 'oldstyle' and 'fancy'
\moderncvcolor{blue}                               % color options 'black', 'blue' (default), 'burgundy', 'green', 'grey', 'orange', 'purple' and 'red'

% adjust the page margins
\usepackage[scale=0.75]{geometry}
\setlength{\footskip}{136.00005pt}                 % depending on the amount of information in the footer, you need to change this value. comment this line out and set it to the size given in the warning

% font loading
% for luatex and xetex, do not use inputenc and fontenc
% see https://tex.stackexchange.com/a/496643
\ifxetexorluatex
  \usepackage{fontspec}
  \usepackage{unicode-math}
  \defaultfontfeatures{Ligatures=TeX}
\else
  \usepackage[T1]{fontenc}
  \usepackage{lmodern}
\fi
\def\justifying{%
  \rightskip=0pt
  \spaceskip=0pt
  \xspaceskip=0pt
  \relax
}

\definecolor{awesome-emerald}{HTML}{00A388}


% document language
\usepackage[ngerman]{babel}  % FIXME: using spanish breaks moderncv

% personal data
\name{Jan}{Wille}
\address{Petermannstr. 12}{30455 Hannover}{}
\email{mail@janwille.de}
\phone{+49~1520~2056206}

\recipient{Personalabteilung}{Sennheiser electronic SE \& Co. KG\\Am Labor 1\\30900 Wedemark}


\begin{document}

\subject{\color{awesome-emerald}{Bewerbung als Software Engineer Embedded (Stellenangebotsnummer: 2500003M)}}
\opening{Sehr geehrte Damen und Herren,}
\makelettertitle
\justifying

die Entwicklung innovativer Embedded-Lösungen begeistert mich seit meinem dualen Studium bei Siemens
und bewegte mich bereits nach Abschluss von diesem innerhalb von Siemens in die Forschung und
Entwicklung zu wechsel. Dort konnte ich parallel zu meinem Masterstudium mein technisches Know-how
im Bereich Firmware-Entwicklung, Cyber-Security und DevOps-Prozesse kontinuierlich vertiefen. Diese
Fertigkeiten möchte ich künftig als Software Engineer Embedded einbringen.

Während meiner aktuellen Tätigkeit bei Siemens war ich maßgeblich an der Firmware-Entwicklung für
ein Antriebssystem des SINAMICS Familie beteiligt. Dazu gehörten die Implementierung von Features
auf STM32-Controllern, die Konzeptionierung automatisierter Softwaretests im Rahmen von
CI/CD-Prozessen sowie die Mitgestaltung von Security-Konzepten im Einklang mit dem Cyber Resilience
Act. Auch die Modernisierung von Update-Verfahren auf AES-Verschlüsselungsstandards und die
Migration der Codebasis von GitLab auf GitHub Enterprise zählten zu meinen Projekten. Dabei war mir
stets die enge Zusammenarbeit mit dem gesamten Team, sowohl Hardware als auch Software, wichtig.

Mein Studium und meine Berufserfahrung haben mir ein breites Fundament vermittelt: von
Echtzeitsystemen, Softwareentwicklung und Mikrocontrollern bis hin zu Themen wie
Industrie-4.0, Robotik und Sensortechnik. Neben meiner technischen Expertise bringe ich eine
strukturierte und selbstorganisierte Arbeitsweise, Teamfähigkeit sowie die Fähigkeit zur klaren
Kommunikation in Deutsch und Englisch (C1) mit.

Besonders begeistert mich an Sennheiser die Möglichkeit, in einem globalen Familienunternehmen mit
starkem Innovationsgeist mitzuwirken und Produkte zu entwickeln, die Menschen weltweit durch
einzigartigen Klang verbinden. Ich bin überzeugt, dass ich mit meiner Erfahrung und meiner
Begeisterung für hochwertige Embedded-Systeme einen wertvollen Beitrag zum Fortschritt Ihres
Unternehmens leisten kann.

Sehr gerne möchte ich Sie in einem persönlichen Gespräch von meiner Motivation und Qualifikation
überzeugen und freue mich auf Ihre Einladung.

\vspace{2\baselineskip}
\enlargethispage{2\baselineskip}
\closing{mit freundlichen Grüßen,}
\makeletterclosing

\end{document}

